\section*{Limitations}
\label{sec:limitations}

\paragraph{The aggregate result is grader-dependent, and we report the weaker side.} Our two judges
agree on the sign and significance of every individual behaviour change and disagree on whether
those changes net out to fewer violations (Table~\ref{tab:substitution}). We treat the held-out
judge as decisive because the primary oracle was the training reward, but this is a judgement call,
not a measurement: a third grader could break the tie either way. What is \emph{not} in question is
the substitution itself.

\paragraph{One optimiser.} The look-ahead arm runs a full ten iterations only for the
preference-tree optimiser. The group-relative arm at $K{=}5$ ran five iterations before being
stopped --- which matched its $K{=}0$ counterpart's \emph{budget} to within 3\% (27.1 vs 27.9
GPU-hours) rather than its iteration count. A $K \times$ optimiser contrast therefore does exist
over iterations 0--5, and we exclude it here by scope rather than for want of data: this paper's
claims are stated at iteration~10, outside that window. Whether the substitution reproduces under
a different loss family across a full ten iterations is open, and remains the single most valuable
remaining experiment.

\paragraph{Single run per arm.} Each arm is one training run. The 96 personas give precision
\emph{within} a run --- which is what every paired contrast here uses --- but there is no
between-run variance estimate, so we cannot separate the intervention from run-to-run variability
in where a given seed's policy lands. This is the limitation we would most like to remove and the
most expensive to remove.

\paragraph{No human coding.} Every behaviour count in this paper is produced by an LLM asked to
apply a coding scheme. We mitigate with a second judge from a different model family that never
participated in training, and we report only effects that survive it --- but no MI-trained human
has coded any of these transcripts, and cross-judge agreement is not the same as validity. Both
judges could share a systematic misreading of what counts as over-praise or as advice. This
matters more than usual here, because our headline is precisely a claim about the \emph{relative
weight} of two violation types.

\paragraph{Ten iterations is long enough to see prevention, not long enough to see a ceiling.} The
$K{=}5$ arm's over-praise is still creeping upward at the endpoint ($0.47 \to 0.63$ acts per
session over iterations 8--10) while the $K{=}0$ arm's roughly doubles over the same stretch. We
can say the channel stayed closed for as long as we ran it, and that the divergence is widening
rather than converging. We cannot say the $K{=}5$ policy would never get there.

\paragraph{The client is a language model.} Change talk, advice-receptivity and the plausibility of
a therapist turn are all properties of a simulated client that shares an architecture and a
training corpus with the grader. A different client could easily make advice-giving a worse route
than flattery, or vice versa, and the specific substitution we observe should not be expected to
transfer. The structural claim --- that closing one route reallocated rather than reduced --- is
what we would expect to transfer, and we have tested it once.

\paragraph{Unseen personas.} The same 96 personas are used for training generation and for
evaluation. There is no held-out-persona generalisation test, so all effects are within the
persona distribution the policy was trained against.

\paragraph{Aggregate choice.} Our aggregate is the instrument's own session total over six coded
behaviours, which weights every violation equally. A clinician might reject that: over-praise and
advice-without-permission are not equally bad. A severity-weighted aggregate would be preferable,
and it is exactly the quantity our graders fail to agree on --- the primary oracle rates the
$K{=}5$ mix more severe at iteration~8 ($\dz{}=-0.59$) while the held-out judge does not
($\dz{}=-0.11$, $n.s.$), and neither effect persists to the endpoint.
